\begin{pastexam}[label={px:gw3-106-ds-03}]{關務三等資訊處理 106 年 第三題}
（申論；25 分）

一個工廠有 $n$ 台機器 M$_1$,M$_2$, $\ldots$,M$_n$ 及 $k$ 份工作 J$_1$, J$_2$, $\ldots$, J$_k$，每份工作都有其所需的執行時間
T（J$_1$），T（J$_2$），$\ldots$,T（J$_k$）。每一台機器一次只能執行一份工作，每份工作只能交給一台機器執行，$n$ 台機器
可同時執行 $n$ 份不同的工作。
\begin{enumerate}[leftmargin=2.6em, labelsep=0.4em]
\item[㈠] 請設計一個 Greedy（貪婪）的演算法，來解決工作排程的問題，使得完成 $k$ 份工作的時間最短。（15 分）
\item[㈡] 此 Greedy 演算法適合使用何種資料結構來完成？（5 分）
\item[㈢] 此 Greedy 演算法的解法是否能保證為最佳解？請舉例說明。（5 分）
\end{enumerate}
\end{pastexam}

\begin{pastsoln}
問題是：把 $k$ 份工作分給 $n$ 台機器，每台機器的負載是分到的工作時間總和，全部完成的時間是最大的負載
（makespan），要讓它最小。 工作不能切開。

\textbf{㈠ 貪心演算法：最長的工作先排}（Longest Processing Time first，LPT）。 先把工作依執行時間由長到短排序，
再依序把每份工作交給「目前負載最小（最早空下來）」的機器。 直覺是：長的工作先放，最後剩下的短工作拿來填平各台
機器的差距。
\begin{pseudo}{LPT 排程}
\Function{LPT-Schedule}{$T, k, n$} \Comment{$T[1..k]$ 是工作時間，$n$ 台機器}
    \State 把工作依 $T$ 由大到小排序
    \State 建立最小堆積 $Q$，放入 $n$ 個 $(0, j)$，$j = 1..n$ \Comment{（負載，機器編號）}
    \For{每份工作 $J$，依排序後的順序}
        \State $(\ell, j) \gets$ \Call{Extract-Min}{$Q$} \Comment{目前負載最小的機器}
        \State 把 $J$ 指派給機器 $M_j$
        \State \Call{Insert}{$Q, (\ell + T(J), j)$}
    \EndFor
    \State \Return 所有機器負載的最大值 \Comment{完成時間}
\EndFunction
\end{pseudo}
時間：排序 $O(k \log k)$，$k$ 次堆積操作各 $O(\log n)$，共 $O(k \log k + k \log n)$。

例：$n = 2$，工作時間 $3, 3, 2, 2, 2$（已由大到小）。 依序：$3 \to M_1$（負載 $3, 0$）；$3 \to M_2$（$3, 3$）；
$2 \to M_1$（$5, 3$，平手取編號小的）；$2 \to M_2$（$5, 5$）；$2 \to M_1$（$7, 5$）。 完成時間 $7$。

\textbf{㈡ 資料結構}：每一步都要找「目前負載最小的機器」，再把它的負載加大放回，正是最小優先佇列的操作，用以
負載為鍵的\textbf{最小堆積}（資結書第~8.3~節），每次 $O(\log n)$。 工作的順序則用排序好的陣列（或以時間為鍵的最大
堆積，每次取最長的）。 若改用陣列每次掃一遍找最小負載，每步 $O(n)$，總共 $O(kn)$。

\textbf{㈢ 不保證最佳}。 上面的例子就是反例：把兩份 $3$ 給 $M_1$、三份 $2$ 給 $M_2$，兩台都是 $6$，完成時間 $6 < 7$。
LPT 把兩份 $3$ 分到不同機器，看起來平均，卻讓三份 $2$ 沒辦法平分。 這不是 LPT 選得不好而已：多台機器的排程問題
（$P \| C_{\max}$）是 NP 困難的，$n = 2$ 時就包含了「把一組數分成和相等的兩堆」的分割問題（資結書第~28~章），所以
目前沒有任何多項式時間的貪心能保證最佳。 LPT 的品質有保證：它的完成時間不超過最佳解的
$\frac{4}{3} - \frac{1}{3n}$ 倍（Graham 1969）；本例 $7 \le (\frac{4}{3} - \frac{1}{6}) \cdot 6 = 7$，剛好達到這個上界。 不排序、
依原順序指派的版本（list scheduling）只保證 $2 - \frac{1}{n}$ 倍（資結書第~29~章的近似演算法觀念）。
\end{pastsoln}
